\documentclass[10pt,a4paper]{amsart}
\usepackage[margin=19mm]{geometry}
\usepackage{amsmath,amssymb,mathtools}
\usepackage[scr=rsfs,cal=euler]{mathalfa}
\usepackage{xfrac}
\usepackage[unicode,hidelinks,bookmarks=false]{hyperref}
\newcommand{\ie}{i.e.\ }
\newcommand{\cf}{cf.\ }
\newcommand{\resp}{respectively\ }
\newcommand{\Subsection}{\subsection}
\providecommand{\nobreakdash}{\nobreak\hbox{-}}
\setlength{\emergencystretch}{2em}
\multlinegap0pt
\usepackage{mdframed}
\usepackage{mdframed}
\usepackage{mdframed}
\usepackage{mdframed}
\usepackage{xcolor}
\usepackage{float}
\usepackage{amssymb}
\usepackage{mathtools}
\usepackage{bm}
\usepackage[shortlabels]{enumitem}
\usepackage{dsfont}
\usepackage{tikz}
\usepackage{stackengine}
\usepackage{mdframed}
\newtheorem{assumption}{Assumption}
\newtheorem{lemma}{Lemma}
\newcommand{\Abar}{\bar{A}}
\newcommand{\E}{\operatorname{\mathbb{E}}}
\newcommand{\Fcal}{\mathcal{F}}
\newcommand{\Gcal}{\mathcal{G}}
\newcommand{\Ocal}{\mathcal{O}}
\newcommand{\R}{\mathbb{R}}				% Real line R
\newcommand{\T}{\top}
\newcommand{\Tcal}{\mathcal{T}}
\newcommand{\bgl}{{\big |}}
\newcommand{\cN}{\mathcal{N}}
\newcommand{\li}{\left<}
\newcommand{\nn}{\nonumber}
\newcommand{\norm}[1]{\|{#1} \|}
\newcommand{\normbig}[1]{\left\|{#1} \right\|}
\newcommand{\ri}{\right>}
\newcommand{\tf}[1]{\left\|{#1}\right\|_{F}}
\newcommand{\tn}[1]{\left\|{#1}\right\|_{\ell_2}}
\newcommand{\xtil}{\hat{x}}
\newcommand{\zbar}{\bar{z}}
\title{Two-Layer Linear Auto-Regressive Models Estimate Latent States}
\author{Yahya Sattar, Sunmook Choi, Leo Maynard-Zhang, Yassir Jedra, Maryam Fazel, Sarah Dean}
\date{Source: arXiv:2606.12691v2 (CC BY 4.0)}
\begin{document}
\maketitle

\noindent\emph{Source and licence.} Two-Layer Linear Auto-Regressive Models Estimate Latent States. Version arXiv:2606.12691v2 (CC BY 4.0), distributed by arXiv under the Creative Commons Attribution 4.0 International licence, http://creativecommons.org/licenses/by/4.0/, as recorded in the arXiv licence metadata for this exact version and shown on the abstract page https://arxiv.org/abs/2606.12691v2. Full licence notice: \texttt{licenses/arxiv-2606.12691-CC-BY-4.0.txt}.\par\smallskip

\noindent\textit{Editorial scope.} Counted unit: the article's lemma lemma:self-normalized\_leq\_1 (source0.tex lines 2225--2230). The article's complete original proof of it is delivered with it (source0.tex lines 2231--2270, one \texttt{proof} environment of 40 lines). No mathematics is written, repaired or re-derived: the statement and the proof are the article's own lines, in the article's own notation, and no retained line is substituted or re-flowed.  Retained uncounted as the source-local context the proof needs: lines 1500--1511; lines 1508--1510; lines 1655--1657; lines 2159--2161.  Nothing was omitted from any retained span, so the package carries no omission record.  No retained line contains a citation, so the wrapper carries no bibliography and no bibliography entry.  No retained line contains a figure environment, an image-inclusion command or drawing code, so no image is delivered and the package declares no asset dependency.  Editorial formatting only, in addition: an AMS \texttt{amsart} preamble that restates the article's own declarations and macros used by the retained lines, namely the environments \texttt{assumption}, \texttt{lemma} on separate counters, together with the standard packages those lines need.  The retained environments print in this document as Assumption 1, Lemma 1 in order of appearance, whereas the article prints the same items with the numbering of its own full document.  The complete native source of the article as distributed in the arXiv e-print for this exact version is bundled unchanged as \texttt{sources/202638/source0.tex}.
Consider a partially observed linear dynamical system with the following state-space representation: for all $t \ge 0$, 
\begin{equation}
\begin{aligned}\label{eqn:linear sys}
	x_{t+1} &= A x_t + Bu_t + w_{t},\\
	y_t &= C x_t + v_t,
\end{aligned}
\end{equation}
where at time $t$, $x_t \in \R^n$ represents the latent state, $u_t \in \R^p$ is the observed control input, $y_t \in \R^m$ is the measured output, $w_t \in \R^n$ is the unobserved process noise, and $v_t \in \R^m$ is the unobserved measurement noise. 
\begin{assumption}\label{assump:Gaussian main}
     During the training period, the noise processes $\lbrace w_t \rbrace_{t\ge0}$, $\lbrace v_t \rbrace_{t\ge 0}$ and the excitations $\lbrace u_t \rbrace_{t\ge 0}$ are sequences of independent, zero-mean, Gaussian random vectors with covariance $\Sigma_w \succ 0$, $\Sigma_v \succ 0$, and $\Sigma_u \succ 0$, respectively.
\end{assumption}
Note that we do not assume knowledge of the matrices $A,B, C$ or the noise covariances $\Sigma_w, \Sigma_v$. Given only a single trajectory of input-output samples $\lbrace (u_t, y_t) \rbrace_{t=0}^{T+H}$ from the system \eqref{eqn:linear sys} satisfying Assumption~\ref{assump:Gaussian main}, our goal is to directly learn the optimal filter, without explicitly learning the system parameters $A,B,C$, and $\Sigma_w, \Sigma_v$.  

\begin{assumption}
     During the training period, the noise processes $\lbrace w_t \rbrace_{t\ge0}$, $\lbrace v_t \rbrace_{t\ge 0}$ and the excitations $\lbrace u_t \rbrace_{t\ge 0}$ are sequences of independent, zero-mean, Gaussian random vectors with covariance $\Sigma_w \succ 0$, $\Sigma_v \succ 0$, and $\Sigma_u \succ 0$, respectively.
\end{assumption}

\begin{assumption} \label{assump:System main}
    (a) The system \eqref{eqn:linear sys} is non-explosive, i.e., $\rho(A) \leq 1$; (b) The pair $(A, \Sigma_w^{1/2})$ is stabilizable, and the pair $(A,C)$ is observable. 
\end{assumption}

\begin{align}
    \zeta_t := \Ocal \Abar^L \xtil_{t-L} + \xi_t, \qquad \xi_t := \Tcal_u u_{t:t+H-2} + \Tcal_e e_{t:t+H-1}, \label{eqn:error definition}
\end{align}

\begin{lemma}[Self-normalized offset bound]\label{lemma:self-normalized_leq_1}
     Suppose Assumptions~\ref{assump:System main}, \ref{assump:Gaussian main} hold. Then, for any $\Theta \in \cN(\bar\Gcal, \varepsilon, \tf{\cdot})$, and $\lambda \in \left[0, 1/\left(18H\norm{\Tcal_u(\Sigma_u \otimes I)\Tcal_u^\T + \Tcal_e(\Sigma_e \otimes I)\Tcal_e^\T}\right)\right]$, we have
    \begin{align}
        \E\left[\exp\left(\lambda \sum_{t=1}^{T} \left(6 \li \xi_t, \Theta \zbar_t \ri - \tn{\Theta\zbar_t}^2 \right) \right)\right] \leq 1.
    \end{align}
\end{lemma}

\begin{proof}
Recall from \eqref{eqn:error definition} that $\xi_t = \Tcal_u u_{t:t+H-2} + \Tcal_e e_{t:t+H-1}$. For the ease of notation, 
%set $K := L+H$, and 
and without loss of generality suppose $N := T/H$ is an integer. Then we have
\begin{align}
    \sum_{t=1}^{T} 6 \li \xi_t, \Theta \zbar_t \ri - \tn{\Theta\zbar_t}^2
    &= \sum_{\tau=0}^{H-1} \sum_{i=0}^{N-1} 6 \li \xi_{1 + iH+ \tau}, \Theta \zbar_{1 + iH+ \tau} \ri - \tn{\Theta\zbar_{1 + iH+ \tau}}^2, \nn \\
    &=: \sum_{\tau=0}^{H-1} \sum_{i=0}^{N-1} 6 \li \xi_{\tau,i}, \Theta \zbar_{\tau,i} \ri - \tn{\Theta\zbar_{\tau,i}}^2, \label{eqn:offset_blocking}
\end{align}
where the subscript $(\tau,i)$ denotes the time index $t= iH+\tau+1$. Applying H\"{o}lder's inequality, we have
\begin{align}
     \E\left[\exp\left(\lambda \sum_{t=1}^{T} 6 \li \xi_t, \Theta \zbar_t \ri - \tn{\Theta\zbar_t}^2 \right)\right] &= \E\left[\exp\left(\lambda \sum_{\tau=0}^{H-1} \sum_{i=0}^{N-1} 6 \li \xi_{\tau,i}, \Theta \zbar_{\tau,i} \ri - \tn{\Theta\zbar_{\tau,i}}^2 \right)\right], \nn \\
     &\leq \prod_{\tau=0}^{H-1} \left(\E\left[\exp\left(\lambda H \sum_{i=0}^{N-1} 6 \li \xi_{\tau,i}, \Theta \zbar_{\tau,i} \ri - \tn{\Theta\zbar_{\tau,i}}^2 \right)\right]\right)^{1/H}. \label{eqn:offset_blocking_holder}
\end{align}
To proceed, let $\{\Fcal_t\}_{t \geq -1}$ be an increasing filtration ($\sigma$-algebra) with all randomness up till time $t-1$. Let $(\tau,i) := iH +\tau+1$ denote a time index, parameterized by $\tau \in [0,H-1]$, and $i \in [0, N-1]$. Then, we have
\begin{align}
    &\E\left[\exp\left(\lambda H \sum_{i=0}^{N-1} 6 \li \xi_{\tau,i}, \Theta \zbar_{\tau,i} \ri - \tn{\Theta\zbar_{\tau,i}}^2\right)\right], \nn \\
    &= \E\bigg[\exp\left(\lambda H \sum_{i=0}^{N-2} 6 \li \xi_{\tau,i}, \Theta \zbar_{\tau,i} \ri - \tn{\Theta\zbar_{\tau,i}}^2\right)\nn \\
    &\quad\quad\quad\quad\E\left[\exp\left(\lambda H \left(6 \li \xi_{\tau,N-1}, \Theta \zbar_{\tau,N-1} \ri - \tn{\Theta\zbar_{\tau,N-1}}^2\right) \right) \bgl \Fcal_{\tau,N-1} \right]\bigg], \nn \\
    &\leq \E\left[\exp\left(\lambda H \sum_{i=0}^{N-2} 6 \li \xi_{\tau,i}, \Theta \zbar_{\tau,i} \ri - \tn{\Theta\zbar_{\tau,i}}^2\right)\right], \label{eqn:offset_tower_N-1}
\end{align}
where we obtained \eqref{eqn:offset_tower_N-1} as follows,
\begin{align}
    &\E\left[\exp\left(6\lambda H \li \xi_{\tau,N-1}, \Theta \zbar_{\tau,N-1} \ri - \lambda H \tn{\Theta\zbar_{\tau,N-1}}^2 \right) \bgl \Fcal_{\tau,N-1} \right] \nn \\ 
    &\qquad= \exp\left(- \lambda H \tn{\Theta\zbar_{\tau,N-1}}^2 \right)\E\left[\exp\left(6\lambda H \li \xi_{\tau,N-1}, \Theta \zbar_{\tau,N-1} \ri  \right) \bgl \Fcal_{\tau,N-1} \right], \nn\\
    &\qquad\leq \exp\left(- \lambda H \tn{\Theta\zbar_{\tau,N-1}}^2 \right) \exp\left(18\lambda^2 H^2\left(\norm{\Tcal_u(\Sigma_u \otimes I)\Tcal_u^\T + \Tcal_e(\Sigma_e \otimes I)\Tcal_e^\T}\right) \tn{\Theta\zbar_{\tau,N-1}}^2\right), \nn \\
    &\qquad= \exp\left(\left(18\lambda^2 H^2\left(\normbig{\Tcal_u(\Sigma_u \otimes I)\Tcal_u^\T + \Tcal_e(\Sigma_e \otimes I)\Tcal_e^\T}\right)- \lambda H\right) \tn{\Theta\zbar_{\tau,N-1}}^2\right), \nn \\
    &\qquad \leq 1,
\end{align}
where we obtained the last inequality by choosing $\lambda \in \left[0, 1/\left(18 H\norm{\Tcal_u(\Sigma_u \otimes I)\Tcal_u^\T + \Tcal_e(\Sigma_e \otimes I)\Tcal_e^\T}\right)\right]$. Repeating the same argument by conditioning on the filtration $\Fcal_{\tau,N-2}, \Fcal_{\tau,N-3}, \dots, \Fcal_{\tau,0}$ in \eqref{eqn:offset_tower_N-1}, we find that
\begin{align}
    \E\left[\exp\left(\lambda H \sum_{i=0}^{N-1} 6 \li \xi_{\tau,i}, \Theta \zbar_{\tau,i} \ri - \tn{\Theta\zbar_{\tau,i}}^2\right)\right] \leq 1, \label{eqn:offset_bound_single_tau}
\end{align}
for $\lambda \in \left[0, 1/\left(18 H\norm{\Tcal_u(\Sigma_u \otimes I)\Tcal_u^\T + \Tcal_e(\Sigma_e \otimes I)\Tcal_e^\T}\right)\right]$. Combining \eqref{eqn:offset_bound_single_tau} with \eqref{eqn:offset_blocking_holder}, we have
\begin{align}
     \E\left[\exp\left(\lambda \sum_{t=1}^{T} 6 \li \xi_t, \Theta \zbar_t \ri - \tn{\Theta\zbar_t}^2 \right)\right] &\leq \prod_{\tau=0}^{H-1} \left(\E\left[\exp\left(\lambda H \sum_{i=0}^{N-1} 6 \li \xi_{\tau,i}, \Theta \zbar_{\tau,i} \ri - \tn{\Theta\zbar_{\tau,i}}^2 \right)\right]\right)^{1/H} \nn \\
     &\leq 1.\label{eqn:offset_blocking_final}
\end{align}
This completes the proof.
\end{proof}

\end{document}
